\section*{PHỤ LỤC}
\phantomsection\addcontentsline{toc}{section}{\numberline{} PHỤ LỤC}

\subsection*{A. Mã nguồn topology.py (Trích đoạn chính)}
\begin{lstlisting}[style=customcode, language=Python, caption={Lớp FRRouter -- Thiết bị định tuyến trên Mininet.}]
class FRRouter(Node):
    def config(self, **params):
        super(FRRouter, self).config(**params)
        self.cmd('sysctl -w net.ipv6.conf.all.forwarding=1')
        self.cmd('sysctl -w net.ipv6.fib_multipath_hash_policy=1')
        self.cmd('sysctl -w net.ipv6.conf.all.accept_dad=0')
        confDir = f'/tmp/{self.name}'
        self.cmd(f'rm -rf {confDir} && mkdir -p {confDir}')
        # Khoi chay Zebra va OSPF6D daemon
        self.cmd(f'/usr/lib/frr/zebra -d -f {confDir}/zebra.conf')
        self.cmd(f'/usr/lib/frr/ospf6d -d -f {confDir}/ospf6d.conf')
\end{lstlisting}

\subsection*{B. Mã nguồn microsegment.sh (Trích đoạn QT1 + QT5)}
\begin{lstlisting}[style=customcode, language=bash, caption={Zero Trust: Cách ly Web và Bảo vệ Database.}]
# [QT1] Web Isolation: Chan Web1 <-> Web2
for web in web_server1 web_server2; do
    ip netns exec $web ip6tables -A INPUT -s fd00:10::/64 -j DROP
done
# [QT5] Web -> DB: ALLOW TCP 3306
ip netns exec s5 ip6tables -A FORWARD -s fd00:10::/64 \
    -d fd00:30::/64 -p tcp --dport 3306 -j ACCEPT
\end{lstlisting}

\subsection*{C. Mã nguồn nat\_setup.sh (Cấu hình Tayga NAT64)}
\begin{lstlisting}[style=customcode, language=bash, caption={Cấu hình NAT64 trên Border Router R1.}]
# Tao giao dien TUN nat64
ip netns exec r1 tayga -c /tmp/r1_tayga/tayga.conf --mktun
ip netns exec r1 ip link set nat64 up
# Routing
ip netns exec r1 ip route add 192.168.255.0/24 dev nat64
ip netns exec r1 ip -6 route add 64:ff9b::/96 dev nat64
# SNAT MASQUERADE
ip netns exec r1 iptables -t nat -A POSTROUTING \
    -s 192.168.255.0/24 -o r1-eth1 -j MASQUERADE
\end{lstlisting}
